\newpage

\section{\huge \textcolor{blue}{一阶自治方程}}
 
\subsection{\LARGE 一阶自治方程}
先给出一阶自治方程的形式:
{\Large
\begin{equation}
\frac{dy}{dt} = f(y)
\end{equation}
}
可看出自治方程的特征就是方程右侧没有$t$, 若用分离变量法可正常求解
但在此，探求在不解方程的情况下判断相关性质:
\begin{enumerate}
    \item 方向场
    \item 临界点
\end{enumerate}

\subsubsection{\large 方向场}
先任给一个自治方程作图如下:
\begin{figure}[htbp]
    \centering
\includegraphics[width=0.6\textwidth]{C:/Users/37/Desktop/Notes on Common Calculus Equations/Notes and pictures/cwffc5.1.1.png}
    \caption{自治方程模拟方向场}
    \label{fig:cwffc5.1.1}
\end{figure}\\
显然, 发现自治方程的积分曲线对平移保持不变, 所以解形状相同.
\\
\\
\subsubsection{\large 临界点}
关于临界点下面为3条等价式:
\begin{enumerate}
    \item $y_0$ 是临界点
    \item $f(y_0) = 0$
    \item $y = y_0$
\end{enumerate}
显然如图:
\begin{figure}[htbp]
    \centering
\includegraphics[width=0.6\textwidth]{C:/Users/37/Desktop/Notes on Common Calculus Equations/Notes and pictures/cwffc5.1.2.png}
    \caption{临界点}
    \label{fig:cwffc5.1.2}
\end{figure}\\
所有积分曲线只会存在临界线两侧, 之间, 
不会与临界线相交\\
下试求临界点:\\
记 $y = $银行存款, $r = $利率(连续), $w = $盗用率(连续)
\[ \frac{dy}{dt} = r y - w \]
先求临界点:
令 $r y - w = 0$
\[ \therefore y = \frac{w}{r} \]
\\
\\
\\作图:\\
\begin{figure}[htbp]
    \centering
\includegraphics[width=0.4\textwidth]{C:/Users/37/Desktop/Notes on Common Calculus Equations/Notes and pictures/cwffc5.1.3.png}
    \caption{临界点}
    \label{fig:cwffc5.1.3}
\end{figure}\\
下分析积分曲线变化:
\begin{figure}[htbp]
    \centering
\includegraphics[width=0.4\textwidth]{C:/Users/37/Desktop/Notes on Common Calculus Equations/Notes and pictures/cwffc5.1.4.png}
    \caption{积分曲线增长变化}
    \label{fig:cwffc5.1.4}
\end{figure}\\
据此些分析大概积分曲线如下:
\begin{figure}[htbp]
    \centering
\includegraphics[width=0.5\textwidth]{C:/Users/37/Desktop/Notes on Common Calculus Equations/Notes and pictures/cwffc5.1.5.png}
    \caption{积分曲线}
    \label{fig:cwffc5.1.5}
\end{figure}\\

\subsubsection{\large 逻辑斯蒂方程}
逻辑斯蒂方程(人口模型):
{\Large
\begin{equation}
\frac{dy}{dt} = k y
\end{equation}
}
其中记$y$为人口数量, $k$为增长率, 若k
为常数, 即为简单人口模型,
 而逻辑斯蒂增长表示, 假设k为常数不合理
 (因资源有限, 换言之, 人口增长率会随人口数量增加而降低)\\
 下从简单的k探讨:\\
 设: $k = a - b y$
 \[ \frac{dy}{dt} = (a - b y) y \]
 临界点:
 \[ (a - b y) y = 0 \Rightarrow y = 0 \text{ 或 } y = \frac{a}{b} \]
 \begin{figure}[htbp]
    \centering
\includegraphics[width=0.3\textwidth]{C:/Users/37/Desktop/Notes on Common Calculus Equations/Notes and pictures/cwffc5.1.6.png}
    \caption{积分曲线增长变化}
    \label{fig:cwffc5.1.6}
\end{figure}

 \begin{figure}[htbp]
    \centering
\includegraphics[width=0.5\textwidth]{C:/Users/37/Desktop/Notes on Common Calculus Equations/Notes and pictures/cwffc5.1.7.png}
    \caption{积分曲线}
    \label{fig:cwffc5.1.7}
\end{figure}

由图 $y = 0 $ 为非稳定临界点, $y = \frac{a}{b}$ 为稳定临界点
, 若积分曲线在一侧聚拢, 一侧发散, 则称半稳定点\\
下在逻辑斯蒂方程基础上, 加入收割概念(收割按照定值, 应用在养殖场)
{\Large
\begin{equation}
\frac{dy}{dt} = (a - b y) y - h
\end{equation}
}
 \begin{figure}[htbp]
    \centering
\includegraphics[width=0.4\textwidth]{C:/Users/37/Desktop/Notes on Common Calculus Equations/Notes and pictures/cwffc5.1.8.png}
    \caption{积分曲线增长变化}
    \label{fig:cwffc5.1.8}
\end{figure}

 \begin{figure}[htbp]
    \centering
\includegraphics[width=0.5\textwidth]{C:/Users/37/Desktop/Notes on Common Calculus Equations/Notes and pictures/cwffc5.1.9.png}
    \caption{积分曲线}
    \label{fig:cwffc5.1.9}
\end{figure}

由图可知, $h_m$ 即为最大收割点, 收割不使种群数目低于此, 即可得到最大收益